% Figaro Paper K: After TradeLens — A Permissionless Bonded Replacement
% Industrial / Systems Engineering register; the TradeLens replacement
% worked example at the inter-logistics perimeter.
% Audience: industrial-engineering and systems-engineering scholars,
% supply-chain-coordination researchers, maritime-shipping practitioners,
% trade-finance professionals.
% Preprint, May 2026

\documentclass[11pt,a4paper]{article}

\usepackage[utf8]{inputenc}
\usepackage[T1]{fontenc}
\usepackage{amsmath,amssymb,amsthm}
\usepackage{mathtools}
\usepackage{booktabs}
\usepackage{graphicx}
\usepackage{hyperref}
\usepackage{cleveref}
\usepackage{enumitem}
\usepackage{xcolor}
\usepackage[margin=1in]{geometry}
\usepackage{natbib}
\bibliographystyle{plainnat}

\newtheorem{theorem}{Theorem}[section]
\newtheorem{proposition}[theorem]{Proposition}
\newtheorem{corollary}[theorem]{Corollary}
\theoremstyle{definition}
\newtheorem{definition}[theorem]{Definition}
\newtheorem{example}[theorem]{Example}
\theoremstyle{remark}
\newtheorem{remark}[theorem]{Remark}

\newcommand{\figaro}{\textsc{Figaro}}
\newcommand{\pay}{P}
\newcommand{\cumval}{G}

\title{%
  After TradeLens:\\
  A Permissionless Bonded Replacement%
}

\author{
  Alessandro Daliana
}

\date{May 2026}

\begin{document}

\maketitle

% ════════════════════════════════════════════════════════════════════
\begin{abstract}
TradeLens (2018--2023), the IBM--Maersk container-shipping consortium
that aimed to produce a tamper-proof shared record of every container's
journey, shut down with only one major non-Maersk carrier on board. The
shutdown was not a software failure. It was a governance failure
visible from the consortium's structural shape: a competitor-controlled
platform, requiring participation from carriers who would not ratify a
competitor's gatekeeping over an industry resource, with high
integration cost and no settlement guarantee binding parties to the
arrangement. The failure was structural, predictable from the shape of
the architecture, and not specific to TradeLens. TradeLens is
one species of a broader pattern: institutional apparatus from
the prior era of trade coordination --- consortium governance,
on-chain voting, contract-engineering toolchains --- ported onto
a new technical substrate without changing the underlying
coordination mechanism.

The architectural alternative we develop operates at the same
perimeter TradeLens attempted: inter-logistics coordination among
ocean carriers, ports of loading and discharge, customs authorities,
freight forwarders, NVOCCs, customs agents, trade-finance parties /
LC issuers, marine insurers, and inspection services. We present a
permissionless, ownerless bonded composition built on the Figaro
kernel's two mechanisms (asymmetric bonding plus buyer dominance with
atomic resolution). The composition has no consortium, no foundation,
no admin function, no upgrade path, and no central party who can be
ratified or refused. Each provider is an off-chain real-world asset
(RWA) whose on-chain participation is mediated by a wallet, with a
human or autonomous agent operating the wallet on the asset's behalf.
The \texttt{rootBuyer} at this perimeter is the buyer-of-record at the
relevant tier --- typically the importer-of-record, the
forwarder-of-record, or the consignee depending on the Incoterm and
the operational arrangement; the kernel is agnostic about which
party occupies the buyer slot. Bonded commitments at each commit lock
skin in the game proportional to the cumulative process value; atomic
resolution by the buyer-of-record settles the whole process on receipt
of the goods at the agreed delivery point. Carriers who want to
participate compose; carriers who do not do not, with no consequence
for either decision.

The schemas the composition draws on are largely existing infrastructure
(commerce, geo, handoff, jurisdiction, fulfilment, GHG family,
proximity); two new schemas (container-seal, Incoterms-2020-as-anchor)
extend the protocol's surface in ways that respect the
extension-doctrine. We also treat the bill-of-lading transferability
question through the cancellable-seller / counter-process pattern that
expresses transferability at the protocol layer without weakening the
kernel's invariants.

The paper is in industrial-engineering register: the contribution is
the composition and its operational properties, not novel
mechanism-design results. Why TradeLens failed and what a bonded
replacement at the same perimeter looks like architecturally are
answerable as supply-chain coordination questions.
\end{abstract}

\medskip
\noindent\textbf{Keywords:}
container shipping, supply-chain coordination, bills of lading,
Incoterms, MLETR, transferable records, process modeling

% ════════════════════════════════════════════════════════════════════
\section{Introduction}
\label{sec:intro}

Container shipping is a coordination problem of unusual scale: a
single 20{,}000-TEU vessel carries cargo from thousands of shippers,
moves through multiple ports under multiple jurisdictions, transfers
to multiple consignees, and is documented by a stack of bills of
lading, customs filings, GHG attestations, and operational handoffs
that no single party originates and no single party owns. The
information problem is real --- duplicated paper-based records,
opaque vessel-position data, contested seal-integrity claims --- and
the industry has been searching for a digital solution for two
decades.

TradeLens, the consortium platform jointly developed by IBM and
Maersk between 2018 and 2023, was the most ambitious attempt at
that solution. Its approach was a permissioned blockchain on which
participants would publish operational records that all parties to
a shipment could read. The platform shut down in early 2023 after
failing to attract more than one major non-Maersk carrier (CMA CGM)
into participation; MSC, Hapag-Lloyd, ONE, and Evergreen all
declined. The platform was technically functional; the consortium
was not.

\paragraph{Reading the failure structurally.}
The standard reading of TradeLens's failure focuses on adoption
dynamics: the network-effects argument that any shared-ledger
platform requires a critical-mass coalition before it is useful.
This is not wrong, but it is incomplete. The critical-mass argument
explains why TradeLens \emph{failed to grow}; it does not explain
why the carriers who declined to join made what was, from their
perspective, a rational decision. The carriers' decision was not
about technology adoption. It was about ratifying a competitor's
control of an industry resource, which a competitor-controlled
platform necessarily asks them to do.

The \emph{governance} structure of the platform was the binding
constraint, not the network effects. A competitor-controlled
platform requires the other competitors to accept the gatekeeping
authority of the controlling competitor; the other competitors
have correctly read this as conceding strategic position and have
declined. The network-effects collapse follows from the governance
choice; it does not cause it. A platform with neutral governance
might have faced its own difficulties, but it would not have faced
the specific difficulty TradeLens did.

\paragraph{The architectural alternative.}
The structural fix is not better governance but no governance at
the platform layer at all. A permissionless, ownerless settlement
primitive --- the Figaro kernel's bonded commitment --- removes the
governance question from the architecture. There is no consortium
to join; there is a public protocol that any wallet can compose
with. There is no gatekeeping authority to refuse; carriers
participate by signing commitments and decline by not signing. The
network-effects question becomes more manageable because the
platform-as-veto problem disappears: a competitor's participation
does not require the competitor's consent to a competitor's
gatekeeping.

\paragraph{Wallets, real-world assets, and the perimeter.}
Every party in the bonded composition is a wallet representing an
off-chain real-world asset (RWA) with productive, commercial value
--- an ocean-going vessel, a port terminal's berth and crane
capacity, a customs authority's clearance service, a forwarder's
multi-modal coordination capability, an inspection service's
on-the-ground presence, a marine insurer's underwriting capacity.
The wallet holds the asset's earnings as token balances and points
to the asset's off-chain credentials through NFTs (a vessel's
classification certificate, a port operator's authority licence, a
customs broker's customs licence, an insurer's policy authorisation).
A human or autonomous agent operates the wallet on the asset's
behalf. Public-authority wallets (port authorities, customs
authorities) participate on the same footing as commercial
providers; the kernel does not distinguish a tax from a fee for
service. A wallet is an active counterparty in the composition
iff the underlying asset is alive and its credentials are
current --- the asset / wallet / operator three-layer split is
load-bearing throughout what follows.

\paragraph{The perimeter the paper develops.}
This paper develops the bonded composition at the same perimeter
TradeLens attempted: \emph{inter-logistics coordination} among the
parties that move a container from origin to destination ---
ocean carriers, ports of loading and discharge, customs
authorities, freight forwarders, NVOCCs, customs agents,
trade-finance parties / letter-of-credit issuers, marine insurers,
inspection services, and cold-chain monitors. The
\texttt{rootBuyer} at this perimeter is the buyer-of-record at
the relevant tier (importer-of-record, forwarder-of-record, or
consignee, depending on the Incoterm and operational arrangement);
the kernel is agnostic about which party occupies the buyer slot.
Granularity within the perimeter is a design choice the
contracting parties make, not a kernel constraint --- the same
movement of goods can be composed as a coarser process (a few
commits to aggregator wallets) or a finer one (more commits to
specialised value-adders). The per-edge bonding equilibrium holds
under either composition; the cohort-pressure topology under
buyer-dominance does shift with cardinality, since coarsening
collapses N small bonded edges into one larger one and changes
which co-sellers carry the weakest-link load. The composition is not a TradeLens
replacement in the sense of a drop-in alternative product; it is a
re-architecture of the inter-logistics coordination problem in
which the consortium-platform layer dissolves into a
permissionless settlement primitive.

\paragraph{Paper organization.}
\Cref{sec:tradelens} reviews the TradeLens history and the
structural reading of its failure.
\Cref{sec:primitive} summarizes the settlement primitive in the
register the present paper uses.
\Cref{sec:assembly} develops the inter-logistics composition: the
organizational DAG (origin to destination through the
inter-logistics value-adders), the schema-attested operational
handoffs along it, and the kernel commit structure that backs it.
\Cref{sec:bonds} works out the bond posture and cohort-compensation
mechanics across the resource providers, with a worked example.
\Cref{sec:schemas} treats the schema requirements: most existing
infrastructure, two new schemas (container-seal-v1 and
Incoterms-2020-as-anchor), and the deliberate non-inclusion of
transferable-BoL semantics in the kernel-layer schema set.
\Cref{sec:bol} treats bill-of-lading transferability as a
protocol-layer composition pattern (CancellableSeller wrapper +
counter-process).
\Cref{sec:comparison} compares the architecture to CargoX,
TradeTrust, MLETR, and TradeLens itself.
\Cref{sec:scope} states the scope conditions and what the
architecture deliberately does not address.
\Cref{sec:conclusion} concludes.

% ════════════════════════════════════════════════════════════════════
\section{What TradeLens Was and Why It Failed}
\label{sec:tradelens}

We treat the failure structurally rather than narratively.

\paragraph{What TradeLens was.}
TradeLens was a permissioned blockchain platform for global
container shipping, jointly developed by IBM and Maersk and
operational from 2018 to 2023. The platform's stated goal was a
shared, tamper-proof record of every container's journey ---
bills of lading, customs filings, port handoffs, vessel positions
--- accessible to all parties in a shipment.

The technical architecture was Hyperledger-Fabric-based: a
permissioned ledger where participating carriers, ports, and
customs authorities could publish operational records. The
platform's use-case design rested on the assumption that all
major ocean carriers would join, producing a single ledger that
shippers could rely on as comprehensive cross-carrier reference.
\citet{jensen_etal_2019} and \citet{sarker_etal_2021} document
the platform's use-case design as it stood mid-deployment; both
papers predate the shutdown and frame the platform around
adoption potential rather than around the governance refusal we
read as the binding constraint below. The governance-was-binding
reading is original to this paper; Jensen and Sarker supply the
use-case baseline against which the reading is offered.

\paragraph{The adoption shape.}
By 2022, only CMA CGM (a Maersk competitor) had joined as a major
non-Maersk carrier. MSC (the largest carrier by capacity),
Hapag-Lloyd, Ocean Network Express (ONE), and Evergreen had all
declined to participate. In November 2022, IBM and Maersk
announced the platform's shutdown; production ended in Q1 2023
\citep{lloyd_loadstar_2022}.

\paragraph{Why the carriers declined.}
The standard adoption-dynamics reading explains the shutdown
through a familiar mechanism: shippers want a single ledger
covering 90\%+ of their cargo volume; without MSC and Hapag-Lloyd
the data was incomplete; without complete data, shippers wouldn't
pay for the platform; without revenue, carriers had no
incremental reason to join; the equilibrium was reflexive
non-adoption. This is correct as far as it goes.

The deeper reading is that the non-Maersk carriers correctly
identified a strategic problem with the platform's governance
shape. A platform controlled by Maersk, with IBM as the technical
operator and Maersk as the principal customer, is structurally a
coordination layer whose governance accountability runs to a
single industry participant; other carriers are being asked to
provision their operational data into a layer whose contested
cases will resolve in that participant's favour. Whatever the
platform's formal governance arrangements, the
controlling-competitor dynamic is the binding constraint on
competitors' participation.

The argument is not that Maersk acted in bad faith or that the
formal governance was inadequate. The argument is that the
platform's shape \emph{is} the problem: a competitor-controlled
infrastructure requires other competitors to ratify the
controller's strategic position relative to the industry, and
there is no formal-governance fix that dissolves this
requirement. MSC and Hapag-Lloyd are not joining a Maersk-controlled
platform regardless of how its governance is configured.

\paragraph{The general lesson.}
TradeLens's failure is not specific to IBM, to Maersk, or to
shipping. It is the recurring pattern of consortium-governance
infrastructure: any platform whose governance accountability runs
to a participant of the industry the platform is meant to serve
will face structural participation refusal from that
participant's competitors. The fix is not better consortium
governance; the fix is no consortium at the platform layer.

Several adjacent industries have run the same experiment with
the same outcome: SWIFT and its alternatives in payments,
various consortium attempts at securities settlement,
healthcare-records consortia. The pattern is robust. The
structural alternative is permissionless, ownerless infrastructure
at the platform layer, with whatever industry-specific arrangements
the participants want to compose on top.

% ════════════════════════════════════════════════════════════════════
\section{The Settlement Primitive}
\label{sec:primitive}

The bonded primitive in \texttt{FigaroCore} has two mechanisms.
We summarize the features the present argument uses.

\paragraph{Asymmetric bonding (Mechanism~1).}
For a transaction with payment $\pay > 0$ and cumulative upstream
value $\cumval \geq \pay$, the buyer locks $2\pay$ and the seller
locks $2\cumval$. Cooperation is the unique profile surviving
iterated elimination of weakly dominated strategies. The
mechanism scales to $N$-party processes through progressive
collateralization: each seller bonds against the cumulative value
of all upstream commitments, producing a mesh of independently
secured bilateral edges, each carrying its own equilibrium.

\paragraph{Buyer dominance with atomic resolution (Mechanism~2).}
Only the root buyer can trigger resolution; resolution settles all
active orders in the process simultaneously or not at all. The
all-or-nothing rule induces a weakest-link subgame among sellers,
producing cooperation pressure of magnitude $\pay_i + 2\cumval_i$
on every other seller from each seller's perspective, without
explicit communication, governance, or managerial layer.

\paragraph{What this implies for shipping coordination.}
Mechanism~1 prices each leg's enforcement in proportional bond
posture without requiring any consortium or platform to underwrite
it. Mechanism~2 enforces the multi-leg coordination at resolution
without requiring any party to govern the multi-leg arrangement.
The two mechanisms together supply the coordination function that
TradeLens attempted to supply through a consortium-governed shared
ledger; they supply it without needing the consortium.

% ════════════════════════════════════════════════════════════════════
\section{The Assembly}
\label{sec:assembly}

We now describe the bonded composition that operates at the
TradeLens perimeter. The composition has two layers that must be
kept distinct: the \emph{organizational DAG} (how the container
physically moves from origin to destination through the
inter-logistics value-adders, documented via schema attestations)
and the \emph{kernel commit structure} (who commits to whom
on-chain, governed by the kernel's \texttt{buyer} $=$
\texttt{rootBuyer} rule). The kernel is DAG-blind; the schema
layer is DAG-aware. The two layers co-exist in one process.

\paragraph{The organizational DAG.}
A container shipment is a sequence of operational handoffs from
origin to destination. We sketch a canonical international flow,
with the inter-logistics value-adders made explicit:

\begin{verbatim}
            shipper-of-record (origin country)
                |
                | handoff: goods presented for shipment
                v
            freight forwarder + NVOCC
                |   (forwarder coordinates multimodal
                |    movement; NVOCC issues house BoL)
                v
            origin inland leg (truck or rail to port)
                |
                | handoff: container stuffed and sealed
                v
            inspection service (pre-shipment, optional)
                |
                | container-seal-v1 inspection
                v
            port-of-loading (terminal services)
                |
                | handoff: container loaded; master BoL issued
                v
            ocean carrier
                |
                | vessel-position stream (geo-v1),
                | container-seal-v1 sequence,
                | cold-chain monitor (if reefer)
                v
            port-of-discharge (terminal services)
                |
                | handoff: container discharged
                v
            customs agent + customs authority
                |   (agent files entries on behalf of
                |    importer; authority bonded for
                |    clearance service and tariff)
                v
            destination inland leg
                |
                | fulfilment-v1 (delivery to consignee)
                v
            consignee (importer-of-record / buyer-at-tier)
\end{verbatim}

Marine insurance and trade-finance / letter-of-credit issuance
operate as parallel value-adders bonded against the same process
without sitting on the cargo-flow line --- the marine insurer's
wallet bonds against the goods' value, the LC issuer's wallet
bonds against the payment-on-documents arrangement, and both are
sellers in the bonded composition without owning a physical
handoff. Variations on the structure (intermodal rail substituted
for inland truck, multiple consolidation points, transhipment ports,
etc.) instantiate the same pattern with the operational divisions
shifted to reflect the route. Every handoff along this DAG is
documented via a schema-typed attestation: \texttt{handoff-v1} for
cargo transfers, \texttt{container-seal-v1} for seal-state changes,
\texttt{geo-v1} for vessel-position updates,
\texttt{jurisdiction-v1} for customs clearance,
\texttt{fulfilment-v1} for delivery confirmation. None of these
attestations is a kernel commit; they are evidentiary records
documenting the operational reality of the container's flow.

\paragraph{The kernel commit structure.}
The kernel commits, by contrast, all share one buyer. The kernel
enforces \texttt{buyer} $=$ \texttt{rootBuyer} in every order in
a process; the rootBuyer at this perimeter is the
\emph{buyer-of-record at the relevant tier}. Depending on the
Incoterm and the operational arrangement, that party can be the
importer-of-record, the consignee, a forwarder-of-record acting
on the importer's behalf, or in seller-arranged shipments the
shipper-of-record itself. The kernel is agnostic about which
party occupies the buyer slot. We use \emph{importer-of-record}
as the running illustration; the apparatus is the same with any
of the alternatives.

The commits document who the importer-of-record pays and how
much: importer $\to$ shipper-of-record (for the cargo's invoice
value, if the Incoterm assigns the import side that obligation),
importer $\to$ freight forwarder, importer $\to$ NVOCC, importer
$\to$ origin inland carrier, importer $\to$ inspection service,
importer $\to$ port-of-loading authority, importer $\to$ ocean
carrier, importer $\to$ port-of-discharge authority, importer
$\to$ customs agent, importer $\to$ customs authority (for the
import duty), importer $\to$ destination inland carrier, importer
$\to$ marine insurer (if buyer-procured cover), importer $\to$
LC issuer (if the payment is letter-of-credit-arranged), and any
other resource provider the operational arrangement draws on.
The commits are independent of who hands the container to whom in
the organizational DAG above. The importer alone calls
\texttt{resolveProcess} on receipt at the agreed delivery point,
settling the process atomically.

\paragraph{Granularity is a design choice.}
The same shipment can be composed coarsely (a few commits to
aggregator wallets --- the forwarder absorbs origin inland +
port-of-loading + carrier + port-of-discharge + destination inland
into one bonded commit) or finely (each value-adder a separate
commit, as enumerated above). The kernel sees only commits; it
does not reify any tier-role or aggregation choice. Granularity is
chosen by the contracting parties to suit their coordination
needs. The per-edge bonding equilibrium is invariant to the
choice (each commit's $2\pay$ / $2\cumval$ posture is independent
of how many other commits the process carries); the cohort-pressure
topology induced by atomic resolution is not, since the magnitude
$\pay_i + 2\cumval_i$ on every co-seller depends on cohort
cardinality. This paper develops the fine-grained version because
the value-adder structure is the substantive content of the
TradeLens-perimeter coordination problem; coarse aggregations are
admissible, may be operationally appropriate in many deployments,
and produce a different cohort-pressure topology that the assembly
designer can engineer for.

\paragraph{Off-chain organizational shapes.}
The conventional inter-logistics chain involves named
organizational shapes (shipper-of-record, consignee, forwarder,
NVOCC, customs broker, importer-of-record, customs agent). Under
the bonded composition the kernel does not see these tier-roles;
each tier-role is whichever wallet the contracting parties
designate at agreement signing. Off-chain organizational shapes
that today perform discovery, curation, brand-as-quality-signal,
or coordination services around the inter-logistics layer can
still be filled by wallets representing the assets that provide
those services and operated by whichever human or agent stack
the asset's owner runs --- but those wallets cannot insert
themselves as kernel intermediaries between the buyer-of-record
and the operational sellers, since the kernel forbids any party
other than the rootBuyer from being the buyer of within-process
commits.

\paragraph{Per-commit schema clauses.}
Each kernel commit carries the schema clauses relevant to the
seller's role.

Each commit pairs the standard \texttt{commerce-v1} payment
clause with the schema clauses appropriate to the seller's role.

\begin{itemize}
  \item \emph{Shipper-of-record}: \texttt{handoff-v1} on goods
    presented for shipment, optional \texttt{ghg-protocol-v1}
    disclosure for manufacturing emissions.
  \item \emph{Freight forwarder}: \texttt{handoff-v1} on cargo
    receipt at consolidation; \texttt{container-seal-v1} applied
    at the consolidation point.
  \item \emph{NVOCC}: house bill-of-lading reference (the BoL
    itself is treated as a protocol-layer composition pattern in
    \Cref{sec:bol}).
  \item \emph{Origin and destination inland carriers}:
    \texttt{handoff-v1} at pickup and dropoff;
    \texttt{geo-v1} pickup and dropoff geohashes.
  \item \emph{Inspection service}:
    \texttt{container-seal-v1} inspection event; optional
    evidence-hash pointer to a pre-shipment inspection report.
  \item \emph{Port-of-loading and port-of-discharge authorities}:
    payment for terminal handling and gate access;
    \texttt{handoff-v1} at port-side transfers;
    \texttt{container-seal-v1} inspected-intact attestations at
    each port handoff.
  \item \emph{Ocean carrier}: an \texttt{incoterms-2020-v1}
    clause anchoring the shipment term and named place;
    \texttt{handoff-v1} on transfer to the vessel and at
    discharge; \texttt{container-seal-v1} sealing events through
    the voyage; \texttt{geo-v1} vessel-position attestations on
    the voyage cadence; \texttt{ghg-protocol-v1} disclosure for
    voyage emissions; for reefer cargo, cold-chain attestations.
  \item \emph{Customs agent}: \texttt{jurisdiction-v1}
    attestation recording the customs authority's determination,
    plus payment for entry-filing services.
  \item \emph{Customs authority}: the import-duty payment lands
    in the customs-authority wallet at this commit. The authority
    has two distinct functions, only one of which the kernel
    sees: as a bonded receiver of a sovereign extraction
    (the duty itself, a tax on the cargo) the kernel sees a
    wallet that committed and a payment that flowed; as the
    sovereign decisional authority on whether the cargo clears
    under the relevant tariff schedule, the authority continues
    to operate off-chain. Both functions sit in the same wallet
    operationally; only the first is kernel-visible.
    \Cref{sec:scope} treats the dual function in detail.
  \item \emph{Marine insurer}: payment of the policy premium;
    the insurer's bond is the underwriting capacity at risk
    against the insured cargo value. Insurance composes with
    the bonded primitive without modifying the kernel.
  \item \emph{LC issuer} (if applicable): payment of the LC
    issuance fee; the issuer's bond is the credit capacity
    backing the documentary-payment arrangement.
\end{itemize}

\paragraph{Operational sub-procurement as separate processes.}
Each value-adder's own procurement (the carrier's fuel + crew +
maintenance, the forwarder's IT + staffing, the inspection
service's labour + equipment, the port authority's terminal
operations, etc.) is a separate process under that wallet's
operator as its own \texttt{rootBuyer}. The kernel does not link
those processes to the importer's. Cross-process coordination is
achieved through within-process cohort-pressure on each side and
through off-chain reconciliation, not through cross-process atomic
resolution. This is the standard wallet-as-RWA composition
pattern: each provider wallet sustains its participation in this
market through receipts that cover the asset's off-chain
operating expenses, and runs its own bonded processes for those
expenses against its own counterparties.

% ════════════════════════════════════════════════════════════════════
\section{Bond Posture and Cascade Mechanics}
\label{sec:bonds}

The composition's structural contribution is that each leg's
risk is priced and propagated by the bonding rule of the kernel,
without consortium intermediation.

\paragraph{Bond posture per commit.}
The kernel maintains a single process-wide cumulative-value
accumulator $G$ that grows monotonically as each commit lands
($G \leftarrow G + P_i$). At every commit the importer-of-record
(as \texttt{rootBuyer}) posts $2P_i$ for that specific payment,
and the seller posts $2G_{\text{at commit time}}$ --- the full
cumulative process value at the moment they commit. Bond posture
is asymmetric: the importer's per-commit bond scales with the
immediate payment, while each seller's bond scales with the
cumulative value of every upstream commit in the process. Sellers
committing later post correspondingly larger bonds; the last
seller bonds against the full landed value of the cargo.

\paragraph{A bond-posture worked example.}
Consider a stylized \$1{,}000 cargo invoice (containerised goods
shipped from an exporting country to an importing country with a
$5\%$ import duty assessed on the CIF declared value), with the
importer-of-record as \texttt{rootBuyer}. The table below traces
one illustrative commit sequence.

\begin{table}[h]
\centering
\small
\begin{tabular}{lrrrr}
\toprule
Seller (commit sequence) & $P_i$ & $G$ after commit & Importer bond $2P_i$ & Seller bond $2G$ \\
\midrule
Shipper-of-record (cargo, root) & $\$700$ & $\$700$ & $\$1{,}400$ & $\$1{,}400$ \\
Marine insurer & $\$15$ & $\$715$ & $\$30$ & $\$1{,}430$ \\
Freight forwarder & $\$30$ & $\$745$ & $\$60$ & $\$1{,}490$ \\
NVOCC (house BoL) & $\$15$ & $\$760$ & $\$30$ & $\$1{,}520$ \\
Origin inland carrier & $\$20$ & $\$780$ & $\$40$ & $\$1{,}560$ \\
Inspection service (pre-shipment) & $\$10$ & $\$790$ & $\$20$ & $\$1{,}580$ \\
Port-of-loading authority & $\$25$ & $\$815$ & $\$50$ & $\$1{,}630$ \\
Ocean carrier & $\$60$ & $\$875$ & $\$120$ & $\$1{,}750$ \\
Port-of-discharge authority & $\$25$ & $\$900$ & $\$50$ & $\$1{,}800$ \\
Customs agent & $\$20$ & $\$920$ & $\$40$ & $\$1{,}840$ \\
Customs authority (import duty) & $\$39$ & $\$959$ & $\$78$ & $\$1{,}918$ \\
Destination inland carrier & $\$25$ & $\$984$ & $\$50$ & $\$1{,}968$ \\
LC issuer (documentary credit fee) & $\$16$ & $\$1{,}000$ & $\$32$ & $\$2{,}000$ \\
\midrule
Total bond locked across the process & --- & --- & $\$2{,}000$ & $\$21{,}886$ \\
\bottomrule
\end{tabular}
\caption{Stylized per-commit bond posture for a \$1{,}000 cargo
shipment under one TradeLens-perimeter process with the
importer-of-record as \texttt{rootBuyer}. The import duty
(\$39, rounded from $0.05 \times \$775 = \$38.75$) is $5\%$ of
the declared CIF value (\$775 = \$700 cargo + \$15 marine
insurance + \$60 ocean freight). $G$ is the process-wide
cumulative-value accumulator, monotonically growing with each
commit. The importer's bond at each commit scales with the
immediate payment; each seller's bond scales with the cumulative
process value at the moment that seller commits. The
customs-authority wallet is an RWA-as-wallet on the same footing
as commercial providers --- a public service whose continued
operation depends on receipts covering its operating expenses.
Per-leg payments are illustrative; absolute magnitudes vary with
cargo class, route, tariff schedule, and the operational
arrangement (open-account vs.\ documentary-credit, FOB vs.\ CIF,
buyer-procured vs.\ seller-procured insurance, etc.).}
\label{tab:k-bond-posture}
\end{table}

The asymmetry between importer-side and seller-cohort bonds is
visible: the importer's aggregate locked bond is \$2{,}000
($= 2 \times P_{\text{cargo invoice}}$, the standard buyer
posture), while the seller cohort's aggregate locked bond is
\$21{,}886, nearly eleven times the importer's. The seller cohort's
bond is a large multiple of the cargo invoice because each seller
bonds against the cumulative $G$ that has grown across upstream
commits, and the inter-logistics perimeter has a substantial
number of value-adders. As the composition adds more value-adders
(e.g., transhipment ports, additional inland legs, multiple
inspection services) the seller-side total grows super-linearly
while the importer's side stays at $2 P_{\text{cargo invoice}}$.

\paragraph{Slip propagation as cohort-compensation pressure.}
When a slip occurs --- a damaged container, a port congestion
delay, a customs hold, a documentary discrepancy, a cold-chain
breach --- the kernel's atomic-resolution rule means
\emph{every} seller's bond in the importer's process is at risk:
the importer alone resolves, and the importer can withhold
resolution while the seller cohort negotiates compensation. The
unit whose performance failed bears the larger share of
cohort-internal compensation; sellers whose service was unaffected
contribute less or not at all. Compensation flows directly from
the seller cohort to the importer before any external mechanism
(specialized mechanism contracts, dispute forums, insurance,
regulators) is engaged. The mechanism is buyer dominance plus
atomic resolution under bond pressure: cooperation is the unique
profile surviving iterated elimination of weakly dominated
strategies, induced by atomic resolution operating on the bonded
mesh. The same structural shape produces the joint-liability
dynamic in Grameen-style microfinance, but operating without the
social-substrate prerequisites Grameen requires (no repeated
interaction, no local information, no exogenous punishment
technology). Buyer dominance with atomic resolution is in this
sense a \emph{social mechanism}: it produces social-coordination
behavior (peer pressure, cohort negotiation, burden-sharing,
coverage of a struggling counterpart) endogenously through the
bond architecture, with each provider wallet's continued
participation in its market making the negotiation rationally
compelled rather than socially expected. The off-chain question of
which operational unit caused the slip --- forwarder coordination
failure, terminal damage, ocean-leg damage, customs inspection
delay, documentary discrepancy at the LC stage --- is adjudicated
by external forums with the bonded record as input.

\paragraph{Cohort pressure across the resource provider set.}
The composition exhibits cohort pressure across the seller set in
the importer's process. Each seller's bond is at risk if any
other seller's slip propagates to a resolution-withholding by the
importer. Sellers therefore have direct economic interest in
each other's reliability: a forwarder whose coordination failures
force repeat compensation negotiations loses repeat
importer-of-record business not through reputation effects but
because the cohort's expected bond-loss exposure on shipments
through that forwarder is direct and measurable; a port whose
congestion consistently produces cascading delays does likewise.
The mechanism operates without a managerial layer; the bonded
structure does the work that a consortium-governed shared ledger
was meant to do through information transparency alone.

% ════════════════════════════════════════════════════════════════════
\section{Schemas}
\label{sec:schemas}

The composition draws mostly on existing schema infrastructure.
Two new schemas are required: a container-seal attestation and an
Incoterms-2020 anchor. The bill-of-lading transferability
question is treated separately in \Cref{sec:bol}.

\paragraph{Existing schemas.}
The composition reuses the following.

\begin{itemize}
  \item \texttt{figaro-commerce-v1} --- per-leg payment commitment.
  \item \texttt{figaro-geo-v1} --- pickup / drop-off geohashes and
    vessel-position updates.
  \item \texttt{figaro-handoff-v1} --- each container handoff at
    transfer points.
  \item \texttt{figaro-jurisdiction-v1} --- customs-clearance
    attestation.
  \item \texttt{figaro-courier-process-v1} --- the carrier's
    per-role event log over the multi-stop carrier process.
  \item \texttt{figaro-fulfilment-v1} --- final delivery to the
    consignee.
  \item \texttt{figaro-ghg-protocol-v1} (with sister schemas for
    ISO 14064, PAS 2050, EN 16258) --- per-leg emissions
    disclosure on every transport edge.
  \item \texttt{figaro-proximity-policy-v1} together with
    \texttt{figaro-proximity-proof-v1} --- geofence enforcement
    at port handoffs where required.
\end{itemize}

These schemas are already on chain and bound to validators; no
further protocol-extension work is needed for them.

\paragraph{New schema: \texttt{figaro-container-seal-v1}.}
Container seals are the integrity mechanism for ocean shipping.
A uniquely-numbered seal is applied at the origin and broken only
at customs or at the destination consignee. Multiple parties
(carrier, ports, customs, consignee) need a single canonical
attestation of when the seal was applied, when it was inspected
intact, and if it was breached.

The schema's Layer~A spec carries:

\begin{itemize}
  \item \texttt{containerNumber} (string, ISO 6346 format)
  \item \texttt{sealNumber} (string)
  \item \texttt{event} (enum: applied, inspected\_intact,
    transferred, breached, removed\_by\_customs)
  \item \texttt{inspectorAddress} (address-hex; the wallet
    making the attestation)
  \item \texttt{timestamp} (ISO 8601 datetime)
  \item \texttt{locationGeohash} (string, 8-character precision)
  \item \texttt{evidenceHash} (bytes-hex, optional content hash
    of a photo or inspection report)
\end{itemize}

The schema is the seed of a \emph{container-integrity} family
that could later cover seal-equivalents in other transport modes
(rail container seals, air-freight tamper-evident packaging, etc.)
without reauthoring the family. The schemaId binds permanently
under the \texttt{SchemaRegistry}'s first-write-wins discipline.

\paragraph{New schema: \texttt{figaro-incoterms-2020-v1}.}
Incoterms~2020 (from the International Chamber of Commerce) are
the international trade vocabulary for risk-transfer points and
cost allocation across legs of a shipment: EXW, FCA, CPT, CIP,
DAP, DPU, DDP, FAS, FOB, CFR, CIF
\citep{icc_incoterms_2020, ramberg2010icc}. The vocabulary is
useful as cross-party shared reference, but the traditional
Incoterms semantics come from a context without the bonded
primitive's invariants and may not all map cleanly. We treat
this carefully because it is the canonical worked example of
how traditional commercial vocabulary smuggles in assumptions
the bonded architecture cannot honor.

Three structural mismatches deserve direct treatment.

\emph{Risk transfer is a Figaro-foreign concept.} Incoterms
encode a discretionary risk-transfer point: ``risk transfers from
seller to buyer at point X.'' Under the bonded primitive, risk
follows bond --- the bonded party is at stake until they perform
per the agreed clauses, and there is no separate
\emph{risk-transfer} concept distinct from the bonding rule. A
clean port from Incoterms to bonded clauses requires reinterpreting
``risk transfer at X'' as ``the seller's
\texttt{handoff-v1} attestation due at X discharges the seller's
bond exposure for the upstream segment.''

\emph{Cost-vs-risk separation collapses.} Incoterms separate
``who pays for transport'' from ``who bears risk during
transport.'' Under the bonded primitive, both follow from process
structure: a party who bonds for a sub-process is the buyer of
that sub-process, and risk is what the bonding rule already
prices. Some Incoterms features (CIP/CIF's seller-pays-insurance
feature, for instance) require composition with a parallel
insurance process rather than direct encoding in the
delivery-clause specification.

\emph{Buyer dominance changes the resolution model.} Incoterms
assume good-faith resolution between parties; the bonded
primitive encodes buyer dominance plus mutually-assured
bond-loss as the equilibrium that produces good-faith resolution.
Some Incoterms' implicit dispute-resolution paths --- which
typically rely on local commercial-court adjudication --- continue
to function under the bonded architecture but at the off-chain
forum the parties name in their agreement, not as an internal
property of the term.

The right schema is therefore not ``Incoterm as a traditional
contract clause'' but \emph{Incoterm as a reference to a
Figaro-native delivery-clause specification}. The schema anchors
the term plus the named place; the term-to-delivery-clause
mapping lives in the off-chain ICC publication (the canonical
text) and a runtime term-to-clause table that the assembly's
participants compose from. The validator contract enforces the
anchor (term enum, named place string, place geohash for
cross-reference with \texttt{geo-v1}) but does not encode
per-term semantics in Solidity.

The Layer~A spec carries:

\begin{itemize}
  \item \texttt{term} (enum: EXW, FCA, CPT, CIP, DAP, DPU, DDP,
    FAS, FOB, CFR, CIF)
  \item \texttt{namedPlace} (string; every Incoterm requires a
    named place)
  \item \texttt{placeGeohash} (string, 8-character precision
    for cross-reference with \texttt{figaro-geo-v1})
\end{itemize}

This split lets future Incoterms revisions ship as new schemaIds
(\texttt{figaro-incoterms-2030-v1}, etc.) without mutating prior
anchors. Per-term mapping verification --- which terms map
directly, which require composition with parallel processes
(insurance, customs), and which do not transfer at all --- is
work the schema-author performs against the kernel code before
the schema is registered. Likely outcomes from that verification
are the following.

\emph{EXW, FCA, DAP, DPU, DDP, FAS, FOB} probably map directly,
each becoming a delivery-clause specification: \texttt{handoff-v1}
attestation at the named place, with auxiliary clauses for
customs (DDP) or unloading (DPU) where the term requires them.

\emph{CPT and CFR} map structurally. The ``seller pays carriage''
feature becomes a separate sub-process where the seller is the
buyer of carriage. The ``risk transfers at first carrier''
feature becomes the goods-commit's delivery-clause specification.

\emph{CIP and CIF} map similarly to CPT and CFR plus the
``seller insures for buyer's benefit'' feature. The insurance
assignment may not transfer directly --- it likely requires
composition with a parallel insurance process where the seller
buys insurance from an external insurer naming the buyer as the
beneficiary. The schema-author's verification report identifies
this as a composition requirement rather than as a refusal.

The Incoterms schema is therefore a clean illustration of the
extension-doctrine principle that the protocol anchors shared
references but does not import every contractual assumption that
travels with traditional vocabulary. The vocabulary is useful;
the assumptions are sometimes incompatible; the schema separates
them.

\paragraph{Schemas the assembly does not include.}
The assembly does not include a multi-currency-cross-leg schema
(would break same-unit comparability and the bonding equilibrium),
a force-majeure-escape schema (would re-introduce the kind of
unilateral exit the kernel forbids), or a transferability schema
internal to \texttt{commerce-v1}. The transferability question is
treated separately in \Cref{sec:bol} as a protocol-layer
composition pattern.

% ════════════════════════════════════════════════════════════════════
\section{Bill of Lading Transferability}
\label{sec:bol}

The bill of lading is the carrier's receipt of cargo, the
contract of carriage, and (in the negotiable case) a document of
title. The transferability question --- whether the right to
receive cargo at destination can be transferred mid-shipment from
the original consignee to a new consignee --- is doctrinally
significant because container-shipping commerce often involves
cargo sale in transit, and the canonical legal artifact for that
sale is endorsement of the bill of lading.

The bonded primitive's kernel does not support party substitution
within a single bonded order. Three kernel invariants each
independently rule it out: the buyer is fixed at root commitment
and identical across every sub-order in the process; the parties
to each commitment are fixed at the moment of commit; and the
no-escape-hatches discipline forbids unilateral exit paths from a
bonded state. A literal endorsement of the BoL during shipment
would require any of these invariants to bend, and they do not
bend.

The architectural direction is that BoL transferability is a
\emph{protocol-layer composition}, not a kernel-layer property.
The candidate pattern is a \emph{cancellable-seller wrapper plus
counter-process}: the original buyer, who wants to transfer the
right to the cargo mid-shipment, commits a parallel
\emph{cancellation process} $P_{\text{cancel}}$ in which each
sub-seller is compensated for early exit through a fee owed by
the original buyer (and bonded under the standard kernel rule),
and the new buyer commits a fresh transport process for the
remaining legs.

\paragraph{Scope of the sketch.}
The pattern is a sketch, not a deliverable mechanism. Three
elements lie outside the kernel-discipline guarantee and must be
specified by the parties at agreement signing rather than relied
on as architectural invariants.

\emph{Sub-seller consent under the wrapper.} The wrapper turns
on programmatic acknowledgement of cancellation by the affected
sub-seller under a fee schedule pre-agreed at commit time. The
sub-seller bonded under the wrapper has consented to its terms;
whether that consent at commit time satisfies the
kernel-discipline reading of counterparty consent at the moment
of cancellation is a design-discipline question the parties
resolve in the off-chain agreement, not an architectural
property the sketch supplies.

\emph{Cargo-handoff operationalization.} A fresh transport
process for the remaining legs requires the original carrier to
release cargo to the new carrier. The cargo handoff is a
real-world event the off-chain agreement attaches an attestation
surface to (a \texttt{handoff-v1} attestation between the two
carriers, with the original buyer counter-signing); the present
sketch does not pin a specific attestation schema for the
handoff because the choice is deployment-shaped.

\emph{Cancellation-fee arithmetic.} The fee schedule must be set
high enough that sub-sellers do not face an early-exit option
they will refuse, and low enough that the transfer use case
remains operational. Specific values are agreement-shaped, not
architectural; the present paper sets out the structure of the
arithmetic without prescribing values.

The pattern is therefore not a recommended deployment artifact in
the way the schemas of \Cref{sec:schemas} are. We include the
sketch because the BoL-transferability question is the
load-bearing concern for any container-shipping deployment, and
because the conventional answer (negotiable bills of lading
endorsed mid-transit) is precisely what the kernel's buyer-fixed
invariant rules out. Deployments requiring BoL transferability
can rely on conventional BoL doctrine through the off-chain
agreement, or constrain their use cases to the non-transferable
region the kernel admits directly.

\paragraph{Why MLETR-style transferable-record analysis applies.}
The MLETR analysis distinguishes records emitted by the kernel
itself --- commitment digests, bond deposits, the cumulative-value
accumulator state, resolution events --- from discretionary
schema-validated attestations such as \texttt{container-seal-v1}
events or \texttt{handoff-v1} attestations. The kernel byproducts
are schema-typed at the kernel level and cryptographically
authenticated; the MLETR Article~10--11 analysis applies to them
directly, with the kernel's structural properties (singularity
through the cumulative-value accumulator, control through
buyer-dominance, integrity through the chain's consensus)
discharging the control and singularity tests for transferable
records. The discretionary schema-validated attestations are
evidentiary artifacts \emph{about} the underlying transferable
record (the commitment), not separate transferable records
inheriting MLETR status in their own right. The cancellable-seller
/ counter-process pattern preserves the kernel-byproduct fit
because the transfer is achieved through bilateral signing on
two separate processes rather than through unilateral endorsement
of a single record; the protocol-layer transfer pattern does not
violate either MLETR test.

\paragraph{What the bonded composition does and does not preserve.}
A useful organizing distinction is between order-form (named
consignee, transferred by endorsement to a designated party) and
bearer-form (transferred by physical delivery of the instrument
without endorsement) bills of lading. Order-form endorsement is
partially expressible through the cancellable-seller pattern above
(subject to the underspecification axes). Bearer-form BoLs are
functionally equivalent to a transfer-without-counterparty-consent
pattern that the bonded composition does not support.
Bearer-instrument practices remain common in commodity trading
(oil cargoes, grain shipments) where the bill is treated as a
near-money instrument \citep{goode_gullifer_2017}; non-bearer
order-form practice predominates in container shipping but is not
the only operational pattern. The bearer-form exclusion is an
honest exclusion: parties whose operations depend on
bearer-instrument liquidity should plan around the conventional
regime for that part of their workflow.

% ════════════════════════════════════════════════════════════════════
\section{Comparison with Existing Approaches}
\label{sec:comparison}

The architecture should be located against the existing landscape
of digital-shipping infrastructures.

\paragraph{TradeLens.}
TradeLens was a permissioned consortium ledger; the architecture
proposed here is a permissionless bonded settlement primitive.
The functional comparison is asymmetric: TradeLens supplied
information transparency without settlement; the bonded
architecture supplies settlement (with information transparency
emerging as a byproduct of the bonded record). The governance
comparison is also asymmetric: TradeLens was governed by
IBM--Maersk; the bonded architecture has no governance at the
kernel layer. The carriers who declined TradeLens did so for a
governance reason; the architecture proposed here removes the
reason.

\paragraph{CargoX.}
CargoX is a blockchain-based bill-of-lading platform that
emphasizes transferable-electronic-record functionality and
MLETR compliance \citep{cargox_2023}. CargoX is a focused product
addressing the BoL transferability problem; the bonded
architecture is a more general settlement primitive whose
relationship to CargoX is composable rather than competitive. A
shipper using the bonded architecture for the bilateral
transport coordination could compose CargoX (or its successor) at
the BoL layer, with the bonded record providing the underlying
settlement evidence and the CargoX layer providing the
transferable-instrument semantics. The two are not in
competition; they address adjacent layers of the stack.

\paragraph{TradeTrust.}
TradeTrust, a Singapore-government-backed framework for
electronic transferable records, supplies open-source tools and
standards for issuing and managing electronic BoLs and similar
instruments \citep{tradetrust_2023}. Like CargoX, it is composable
with the bonded architecture rather than competitive with it.
TradeTrust's standards-based approach to MLETR fit produces
artifacts that are interoperable with the bonded composition's
evidence record.

\paragraph{MLETR.}
The UNCITRAL Model Law on Electronic Transferable Records (2017)
is a legal framework rather than a technology platform
\citep{uncitral_mletr_2017}. The kernel-layer artifacts satisfy
Articles~10 and 11 (functional equivalence, control, singularity)
for non-bearer transferable records, with bearer-instrument
semantics noted as an honest exclusion (\Cref{sec:bol}).

\paragraph{The architectural distinction.}
The TradeLens replacement is not a better TradeLens. It is a
re-architecture of the underlying coordination function from
\emph{shared ledger maintained by consortium} to
\emph{bonded settlement primitive composed by participants}. The
shift is structurally substantial: the consortium's role is
absorbed into the kernel's permissionless protocol; the carriers'
participation decision is reduced to whether to compose with the
protocol on individual shipments rather than whether to ratify a
competitor's gatekeeping. The carriers who refused to ratify
TradeLens are not asked to ratify anything in the bonded
architecture; they participate by signing commitments to specific
shipments and decline by not signing. The decision is made at
the level of individual commerce rather than at the level of
industry governance.

% ════════════════════════════════════════════════════════════════════
\section{Scope}
\label{sec:scope}

The argument bounds itself by several scope conditions.

\paragraph{Asset specificity at the port layer.}
Major ports are functional monopolies for the slot, terminal
capacity, and operational infrastructure they provide; carriers
have limited alternative-supplier leverage at congested hubs.
The architecture's competitive properties at the port-supply
layer are constrained accordingly: where the port-of-call is
dictated by the route and the port operator is a near-monopolist,
the buyer-dominance leverage at the sub-procurement is
attenuated. The architecture continues to function in these
cases (the bond posture is well-defined; the resolution path is
clean) but it does not produce price competition in port
services where structural conditions don't admit it.

\paragraph{Customs and sovereign authorities.}
Customs authorities have two distinct functions in this
composition. As bonded sellers of border-clearance service
(charging the importer for handling, examination capacity,
filing-receipt processing, etc.) they are RWA-as-wallet on the
same footing as commercial providers, and their participation in
the bonded process is unproblematic at the kernel layer. As
sovereign decisional authorities on whether specific goods clear
under the relevant tariff schedule, they continue to operate in
their own register, off-chain; the architecture does not
re-engineer customs administration and the kernel does not
encode the sovereign decision. The bonded composition produces a
verifiable evidence record on which customs decisions can operate
and which customs decisions are recorded against; this is an
improvement over the paper-and-database baseline but not a
substitute for the sovereign authority itself.

\paragraph{Bond-currency at the small-shipper end.}
Small shippers without crypto-treasury infrastructure face
operational friction at the bonding step. The constraint is
real but solvable through standard treasury-of-record
arrangements (a forwarder bonds on behalf of small shippers,
with the small shipper paying conventional fees off-chain to the
forwarder for the bonding service). This is a runtime-tier
implementation question rather than a substrate-tier question;
the substrate is denomination-agnostic.

\paragraph{Insurance and reinsurance.}
Marine insurance, P\&I clubs, and the general apparatus of
shipping-risk underwriting continue to operate. The bonded
architecture changes the risk profile that the marine-insurance
markets price (smaller residual after bond allocation;
concentrated risk at the carrier-bond layer; a new class of
bond-default risk at the port-supplier layer) but does not
displace the insurance function. We expect the marine-insurance
markets to develop fluent products around the bonded
architecture as the architecture diffuses.

\paragraph{Network effects and adoption.}
The architecture's adoption dynamics differ from TradeLens's in
that no participant must ratify any other participant's
governance position. A shipper, forwarder, or carrier can
participate on individual shipments without joining anything;
the participation decision is shipment-by-shipment rather than
platform-membership-shaped. We do not model the diffusion
trajectory beyond noting that the structural barrier that
defeated TradeLens (consortium ratification by competitors) is
absent here.

% ════════════════════════════════════════════════════════════════════
\section{Conclusion}
\label{sec:conclusion}

TradeLens failed for governance reasons that no
consortium-governance fix can resolve, and the structural
alternative is permissionless, ownerless infrastructure at the
platform layer rather than a better consortium. The bonded
primitive in \texttt{FigaroCore} supplies the architectural
shape: each leg of a shipment is a bilateral bonded commitment;
progressive collateralization scales the primitive across the
multi-leg fan-out; buyer dominance plus atomic resolution
propagates damage to the unit causing the slip; the schemas the
composition reuses are largely existing infrastructure with two
well-bounded additions.

The composition operates at the same perimeter TradeLens attempted
--- inter-logistics coordination among carriers, ports, customs
authorities, freight forwarders, NVOCCs, customs agents,
trade-finance parties, marine insurers, and inspection services
--- with the buyer-of-record at the relevant tier as the
\texttt{rootBuyer}. Each provider is an RWA-as-wallet whose
participation in this market is sustained by receipts covering
the asset's off-chain operating expenses; public-authority
wallets participate on the same footing as commercial.
Granularity within the perimeter is a kernel-blind design choice;
the fine-grained version developed here makes the value-adder
structure explicit, and coarser aggregations are admissible
where operational arrangements support them.

What the present paper supplies is the composition --- the
organizational DAG, the per-edge bonding mechanics, the
schema set including the two new schemas
(\texttt{container-seal-v1}, \texttt{incoterms-2020-v1}), and
the cancellable-seller / counter-process pattern for BoL
transferability --- that turns the bonded primitive into a
working inter-logistics coordination architecture at the same
perimeter TradeLens attempted and could not sustain.

\section*{Acknowledgments}
The TradeLens worked example originated in a project repository
of agent-coordination scenarios. The present paper formalises
the scenario into an industrial-engineering register.

% ════════════════════════════════════════════════════════════════════
\begin{thebibliography}{99}

\bibitem[Jensen et al.(2019)]{jensen_etal_2019}
Jensen, T., Hedman, J., and Henningsson, S.
\newblock How TradeLens Delivers Business Value With Blockchain
  Technology.
\newblock \emph{MIS Quarterly Executive}, 18(4):221--243, 2019.

\bibitem[Sarker et al.(2021)]{sarker_etal_2021}
Sarker, S., Henningsson, S., Jensen, T., and Hedman, J.
\newblock Use Cases for Blockchain in the Enterprise: TradeLens.
\newblock \emph{MIS Quarterly Executive}, 20(4):257--276, 2021.

\bibitem[Lloyd's Loadstar(2022)]{lloyd_loadstar_2022}
The Loadstar.
\newblock IBM and Maersk to Sink TradeLens Blockchain Platform.
\newblock The Loadstar, November~30, 2022.

\bibitem[ICC(2020)]{icc_incoterms_2020}
International Chamber of Commerce.
\newblock \emph{Incoterms 2020: ICC Rules for the Use of Domestic and
  International Trade Terms}.
\newblock ICC Publication 723E, 2020.

\bibitem[Ramberg(2010)]{ramberg2010icc}
Ramberg, J.
\newblock \emph{ICC Guide to Incoterms 2010}.
\newblock ICC Publication 720E, 2010.

\bibitem[CargoX(2023)]{cargox_2023}
CargoX d.o.o.
\newblock \emph{CargoX Platform: Blockchain Document Transfer
  Solution for Shipping}.
\newblock CargoX product documentation, 2023.

\bibitem[TradeTrust(2023)]{tradetrust_2023}
Government of Singapore.
\newblock \emph{TradeTrust: Open-Source Framework for Electronic
  Transferable Records}.
\newblock TradeTrust technical documentation, 2023.

\bibitem[UNCITRAL(2017)]{uncitral_mletr_2017}
United Nations Commission on International Trade Law.
\newblock \emph{UNCITRAL Model Law on Electronic Transferable
  Records}.
\newblock UN Publication, 2017.

\bibitem[Goode and Gullifer(2017)]{goode_gullifer_2017}
Goode, R. and Gullifer, L.
\newblock \emph{Goode and Gullifer on Legal Problems of Credit and
  Security}, 6th edition.
\newblock Sweet \& Maxwell, London, 2017.

\end{thebibliography}

\end{document}
